% Generated draft: requires mathematical and source review.
\documentclass[11pt]{article}
\usepackage[margin=25mm]{geometry}
\usepackage{fontspec}
\setmainfont{Latin Modern Roman}
\usepackage{amsmath,amssymb,mathtools,unicode-math}
\setmathfont{Latin Modern Math}
\usepackage{enumitem,needspace,xurl,hyperref}
\setlength{\parindent}{0pt}
\setlength{\parskip}{5pt}
\setlength{\emergencystretch}{3em}
\newcommand{\N}{\mathbb{N}}
\newcommand{\Z}{\mathbb{Z}}
\newcommand{\Q}{\mathbb{Q}}
\newcommand{\R}{\mathbb{R}}
\newcommand{\C}{\mathbb{C}}
\newcommand{\F}{\mathbb{F}}
\newcommand{\A}{\mathbb{A}}
\newcommand{\PP}{\mathbb{P}}
\newcommand{\E}{\mathbb{E}}
\newcommand{\eps}{\varepsilon}
\newcommand{\dd}{\,\mathrm{d}}
\begin{document}
\section*{UnsolvedMath ID: 800193}
\textbf{Source problem code:} catalogue entry
\subsection{Definitions and mathematical statement}
Let \(X\) and \(X'\) be hyperk\"ahler manifolds with \(\dim_{\C}X=\dim_{\C}X'\). The supplied excerpt does not define the class of hyperk\"ahler manifolds intended. Under the customary compact irreducible holomorphic symplectic interpretation, such a manifold is a simply connected compact K\"ahler manifold whose space of holomorphic two-forms is generated by an everywhere nondegenerate form. The formulation below uses this interpretation conditionally; its applicability to the intended source convention requires review.

Write \(q_X\) and \(q_{X'}\) for the Beauville--Bogomolov integral quadratic forms on second cohomology. A Hodge isometry is an isomorphism of abelian groups
\[
\lambda:H^2(X,\Z)\xrightarrow{\sim}H^2(X',\Z)
\]
that satisfies
\[
q_{X'}(\lambda(\alpha))=q_X(\alpha)
\qquad\text{for every }\alpha\in H^2(X,\Z),
\]
and whose complexification satisfies
\[
\lambda_{\C}\bigl(H^{p,q}(X)\bigr)=H^{p,q}(X')
\qquad(p,q\geq 0,\ p+q=2).
\]
The identification of the quadratic forms with the Beauville--Bogomolov forms is the standard interpretation of the source's phrase ``Hodge isometry''; the supplied excerpt does not state it explicitly.

\textbf{Question (Beauville, 2010, Question 2).} If such a Hodge isometry exists, what geometric relationship between \(X\) and \(X'\) follows? What additional structure, required to be preserved by \(\lambda\), would suffice to conclude that \(X\) and \(X'\) are biholomorphic?

The additional structure is deliberately unspecified in the supplied question, so this is a request for suitable hypotheses rather than a single stated implication with fixed extra assumptions. The source asks for the existence of an isomorphism between the manifolds; it does not explicitly require that the given \(\lambda\) be induced by that isomorphism.
\subsection{Short English statement}
Beauville asks what can be deduced about two hyperk\"ahler manifolds of equal dimension when their integral second cohomology groups are Hodge isometric, and which additional preserved structure would guarantee that the manifolds are biholomorphic. The intended hyperk\"ahler convention and the extra structure require clarification. The supplied registry labels the question open; this draft does not independently verify that status.
\subsection{Sources}
\begin{itemize}
\item Supplied local UnsolvedMath registry record, permanent ID \(800193\), exact source problem code \texttt{catalogue entry}, titled ``Holomorphic Symplectic Geometry --- Question 2,'' in the AMR Open Problem Lists collection. The record attributes the question to Beauville, with proposed year 2010.
\item Beauville, \emph{Holomorphic Symplectic Geometry: A Problem List} (2010), Question 2. Supplied source URL: \url{https://arxiv.org/abs/1002.4321}. The record reports extraction from arXiv LaTeX and access on 2026-10-16; these are supplied provenance statements, not independently verified access or source checks.
\item The supplied registry status is \texttt{open}; the record was created and updated on 2026-10-16. This is a supplied registry label, not an independently verified assessment of the mathematical literature.
\end{itemize}
\end{document}
